\section{Don't Care Conditions}

\subsection{Definition}
A don't care condition occurs when the output of a logic circuit for a
specific input combination is irrelevant. These inputs are represented by \textit{d}
in Karnaugh-maps. Other representations that can be used are X, $\phi$
or $-$, these are commonly used in truth tables.

\subsection{Why They Arise}

Don't care conditions typically arise in two situations:

\begin{description}
    \item[Invalid input combinations.] In a BCD (Binary Coded Decimal)
    system, only the input combinations \(0000\) to \(1001\) (representing
    digits 0--9) are valid. The remaining six combinations (\(1010\) to
    \(1111\)) can never occur, so their outputs are assigned as don't cares.

    \item[Unused outputs.] In some systems a particular input state is
    reachable, but the resulting output is never used by any other part of
    the circuit, making it irrelevant to the designer.
\end{description}

\subsection{How They Are Used}

During minimisation, don't care entries may be treated as either 0 or 1 —
whichever choice produces the simplest expression. This gives the designer
freedom to form \textbf{larger groupings (implicants)}, which directly
reduces the number of logic gates in the final circuit. A key rule applies:
a grouping may \emph{include} don't cares to enlarge it, but it cannot
consist \emph{entirely} of don't cares, as such a group covers no required
output.

\subsection{Truth Table Example}

Consider the function \(f(A,B,C) = \sum m(1,3,7) + \sum d(2,5)\), shown
in Table~\ref{tab:dontcare}. Minterms 2 and 5 are don't cares and can be
included in groupings during minimisation to yield a simpler final
expression.

\begin{table}[h!]
\centering
\caption{Truth table with don't care conditions}
\label{tab:dontcare}
\begin{tabular}{@{}ccccc@{}}
\toprule
\textbf{Minterm} & \textbf{A} & \textbf{B} & \textbf{C} &
\textbf{Output} \\
\midrule
0 & 0 & 0 & 0 & 0 \\
1 & 0 & 0 & 1 & 1 \\
2 & 0 & 1 & 0 & \(d\) \\
3 & 0 & 1 & 1 & 1 \\
4 & 1 & 0 & 0 & 0 \\
5 & 1 & 0 & 1 & \(d\) \\
6 & 1 & 1 & 0 & 0 \\
7 & 1 & 1 & 1 & 1 \\
\bottomrule
\end{tabular}
\end{table}
